\section[Stage 029]{Stage 029: Tracking-Branch Bounds and Residual Comparison}
\label{stage:029}

\stagefield{Part / anchor}{Part III; primary anchor \resultanchor{MTDC-T6}.}
\claimstatus{\StatusExactClosure{} for the monotonicity and residual comparison on the tracking branch.}
\stagefield{Purpose}{Stage~029 determines whether the first coherent tracking deformation helps or hurts the selected normalization target.}
\stagefield{Inputs}{The tracking functions \(F_{\rm tr}\), \(G_{\rm tr}\) and the constructive inequality \(R_{\rm tr}<1\).}

\stagefield{Verification}{SymPy audit: \StageFile{scripts/moving_throat_pde_stage029_tracking_branch_bounds_sympy_audit.py}; transcript: \StageFile{scripts/output/moving_throat_pde_stage029_tracking_branch_bounds_sympy_audit.txt}.  Mathematica audit: \StageFile{mathematica/moving_throat_pde_stage029_tracking_branch_bounds_mathematica_audit.wl}; transcript: \StageFile{mathematica/output/moving_throat_pde_stage029_tracking_branch_bounds_mathematica_audit.txt}.}

\subsection*{Derivation}

The tracking functions obey
\begin{equation}
\label{eq:app-stage029-monotonicity}
\boxed{
\frac{\partial G_{\rm tr}}{\partial R}<0,
\qquad
\frac{\partial F_{\rm tr}}{\partial R}>0}.
\end{equation}
Because the constructive split--\(U\) branch has \(R_{\rm tr}<1\), one finds
\begin{equation}
\label{eq:app-stage029-inequalities}
G_{\rm tr}>G_{\rm flat},
\qquad
F_{\rm tr}<F_{\rm flat}
\end{equation}
at fixed \((\xi,\delta)\).  The exact residual comparison can be written as
\begin{equation}
\label{eq:app-stage029-residual}
\boxed{
E_{\rm tr}-E_{\rm flat}=F_{\rm flat}-F_{\rm tr}>0}.
\end{equation}
Thus the first coherent local D/N kernel worsens the selected normalization relative to the flat branch.  Support enhancement must compensate this deficit.

\stagefield{Output}{Monotonicity \eqref{eq:app-stage029-monotonicity} and residual theorem \eqref{eq:app-stage029-residual}.}
\stagefield{Downstream use}{Stage~030 introduces support enhancement; Stage~031 proves the compensation theorem.}

